\documentclass[10pt]{article}
\usepackage[T1]{fontenc}
\usepackage[utf8]{inputenc}
\usepackage[french]{babel}
\usepackage{amsmath,amssymb,mathrsfs,mathpazo}
\usepackage{geometry}
\geometry{a4paper,margin=2cm}
\usepackage{graphicx,array,tabularx,tabularray,colortbl,multirow,multicol,booktabs,enumitem,pifont,ulem,fancybox,diagbox,cancel}
\usepackage{pstricks,pstricks-add,pst-plot,pst-tree,pst-node,pst-text,pst-func,pst-math,pst-bspline,pst-3dplot,pst-eucl,pst-solides3d}
\let\pscalculate\undefined
\usepackage{pst-bezier}
\usepackage{tikz,pgfplots,tkz-tab}
\usetikzlibrary{arrows,arrows.meta,patterns,calc,positioning,trees,decorations.markings,shapes.geometric,shapes.misc,shapes.arrows,intersections,decorations.pathmorphing,decorations.pathreplacing,backgrounds,shadows,babel}
\pgfplotsset{compat=1.18}
\usepgfplotslibrary{fillbetween}
\usepackage[np]{numprint}
\usepackage{listings,makecell,pas-tableur,twoopt}
\let\barre\undefined
\usepackage[pstricks]{tablvar}
\definecolor{darkgreen}{RGB}{0,100,0}
\providecommand{\red}{\color{red}}
\providecommand{\blue}{\color{blue}}
\providecommand{\green}{\color{green}}
\providecommand{\black}{\color{black}}
\providecommand{\white}{\color{white}}
\providecommand{\gray}{\color{gray}}
\providecommand{\cyan}{\color{cyan}}
\providecommand{\magenta}{\color{magenta}}
\providecommand{\yellow}{\color{yellow}}
\providecommand{\violet}{\color{violet}}
\providecommand{\orange}{\color{orange}}
\providecommand{\euro}{\texteuro}
\providecommand{\EUR}{\texteuro}
\providecommand{\ieme}{\textsuperscript{e}}
\providecommand{\degres}{\ensuremath{^\circ}}
\providecommand{\degree}{\ensuremath{^\circ}}
\providecommand{\vv}[1]{\overrightarrow{#1}}
\providecommand{\Coord}[1]{#1}
\providecommand{\decosix}{}
\providecommand{\starredbullet}{\ensuremath{\bullet}}
\pagestyle{empty}
\setlength{\parindent}{0pt}
\usepackage{adjustbox}\DeclareUnicodeCharacter{2236}{:}\DeclareUnicodeCharacter{2190}{\ensuremath{\leftarrow}}
\providecommand{\R}{\mathbb{R}}
\providecommand{\N}{\mathbb{N}}
\providecommand{\Z}{\mathbb{Z}}
\providecommand{\Q}{\mathbb{Q}}
\providecommand{\D}{\mathbb{D}}
\providecommand{\C}{\mathbb{C}}
\providecommand{\e}{{\mathrm{e}}}
\providecommand{\vect}[1]{\overrightarrow{#1}}
\providecommand{\vectt}[1]{\overrightarrow{\mathrm{#1}}}
\providecommand{\vvt}[1]{\overrightarrow{\mathrm{#1}}}
\providecommand{\np}[2][]{{#2~#1}}
\providecommand{\vv}[1]{{\overrightarrow{#1}}}
\renewcommand{\d}{{\mathrm{d}}}
\providecommand{\degres}{{^\circ}}
\providecommand{\degre}{{^\circ}}
\providecommand{\blue}{\color{blue}}
\providecommand{\red}{\color{red}}
\providecommand{\enskip}{\quad}
\providecommand{\llbracket}{\lbrack\!\lbrack}
\providecommand{\rrbracket}{\rbrack\!\rbrack}
\providecommand{\euro}{{\text{€}}}
\providecommand{\no}{{\text{n°}}}
\providecommand{\interval}[2]{{[#1;#2]}}
\providecommand{\intervFO}[2]{{[#1;#2[}}
\providecommand{\dotfill}{\ldots}

\providecommand{\ds}{\displaystyle}
\providecommand{\textding}[1]{#1}
\providecommand{\Oij}{\(\mathrm{O};\vec{\imath},\vec{\jmath}\)}
\providecommand{\Oijk}{\(\mathrm{O};\vec{\imath},\vec{\jmath},\vec{k}\)}
\providecommand{\Ouv}{\(\mathrm{O};\vec{u},\vec{v}\)}
\providecommand{\intervFF}[2]{\ensuremath{[#1;#2]}}
\renewcommand{\euro}{\ifmmode\text{\texteuro}\else\texteuro\fi}
\DeclareUnicodeCharacter{00B0}{\ifmmode^{\circ}\else\textdegree\fi}
\DeclareUnicodeCharacter{20AC}{\ifmmode\text{\texteuro}\else\texteuro\fi}
% Énoncés : Éducation nationale. Transcription et corrigés : APMEP, Denis Vergès et contributeurs. Découpage et classement : CoopMaths, Sébastien Lozano, Éric Elter et contributeurs. Adaptation de la présentation : Maths Explorer.
% CC BY-SA : https://coopmaths.fr/www/about/
% https://coopmaths.fr/alea/static/e3c/2020/tex/e3c_2020_00_sujet34_1.tex
\begin{document}
\noindent\textit{APMEP / CoopMaths. CC BY-SA. Présentation adaptée par Maths Explorer.}\par\medskip
\medskip
Cet exercice est un questionnaire à choix multiples. 

Pour chaque question une seule réponse est exacte. 

Une mauvaise réponse ou une absence de réponse n'enlève aucun point. 

La bonne réponse rapporte un point.

Il n'est pas demandé de justification.

\medskip

\begin{enumerate}
\item L'ensemble des solutions de l'inéquation $-3x^2+2x+1>0$, où $x$ est un nombre réel, est :

\begin{tabularx}{\linewidth}{*{2}{X}}
\textbf{A.~~} $\left\{-\dfrac{1}{3}~;~1\right\} $ &\textbf{B.~~} $\emptyset$\\\textbf{C.~~ }$\left]-\dfrac{1}{3}~;~1\right[  $& \textbf{D.~~}$\left]-\infty~;~-\dfrac{1}{3}\right[\cup\left]1~;~+\infty\right[ $.
\end{tabularx}
\medskip

\item Le plan est muni d'un repère $(O;\vv{i},\vv{j})$.

Une équation cartésienne de la droite $(d)$ passant par le point A de coordonnées $(-1~;~5)$ et de vecteur directeur $\vv{v}$ de coordonnées $(3~;~-2)$ est :

\medskip

\begin{tabularx}{\linewidth}{*{4}{X}}
\textbf{A.~~} $-2x+3y+13=0$&\textbf{B.~~}$-2x-3y-13=0 $&\textbf{C.~~}$2x - 3y + 13 = 0 $& \textbf{D.~~} \small $-2x - 3y + 13 = 0$.
\end{tabularx}

\medskip

\item Soit $f$ la fonction définie sur $]-\infty~;~2[ \cup  ]2~;~+\infty[$ par $f(x)=\dfrac{2x+1}{x - 2}$.

La fonction dérivée de $f$ est définie sur $]-\infty~;~2[ \cup  ]2~;~+\infty[$ par :

\medskip

\begin{tabularx}{\linewidth}{*{4}{X}}
\textbf{A.~~} $ f'(x)=\dfrac{5}{(x-2)^2}$ &\textbf{B.~~} $f'(x)= \dfrac{3x-6}{(x-2)^2} $&\textbf{C.~~}$f'(x)=\dfrac{-3}{(x-2)^2} $& \textbf{D.~~} $ f'(x) = \dfrac{-5}{(x-2)^2 }$.
\end{tabularx}

\medskip

\item Pour tout nombre réel $x$, une expression simplifiée de \[\dfrac{(\e^x)^2\times \e^{-x+1}}{\e^{5x}}\] est :

\medskip

\begin{tabularx}{\linewidth}{*{4}{X}}
\textbf{A.~~} $\e^{-4x+1} $ &\textbf{B.~~} $\e^{x^2-6x+1} $&\textbf{C.~~}$ \e^{x^2+4x+1}$& \textbf{D.~~} $\e^{-x^3+x^5-5x}  $.
\end{tabularx}

\medskip

\item La fonction $f$ est définie pour tout $x$ réel par \[f(x)=\e^x(3\e^x-1).\]
La fonction dérivée de $f$ est définie pour tout $x$ réel par :

\medskip

\begin{tabularx}{\linewidth}{*{4}{X}}
\textbf{A.~~} $f'(x)=\e^x(3\e^x) $ &\textbf{B.~~} $f'(x)=6\e^{2x}-\e^x $&\textbf{C.~~}$f'(x)=3\e^{2x}-\e{x}$& \textbf{D.~~} $f'(x)=3(\e^x)^2-1  $.
\end{tabularx}
\end{enumerate}

\vspace{0,5cm}

\end{document}
